
\section{Findings and Discussion} 

We summarize key findings from our empirical study and discuss promising ways for addressing GraphRAG's two key bottlenecks: subgraph-extraction efficiency and LLM generation quality.



%\vspace{-4pt}
\subsection{Summary of Key Findings}
%\begin{itemize}[leftmargin=15pt]
%    \item 
%\vspace{-5pt}
{\bf Finding 1: GraphRAG Trade-off Structure.} GraphRAG entails multi-aspect trade-offs among quality (Hits@1), efficiency (runtime), and cost (token and GPU usage). This trade-off spans two modules (subgraph-extraction vs. path-retrieval), two method types (structure-based vs. semantic-augmented), and two types of semantic models (EEMs vs. LLMs).

    %\item 
    \textbf{Finding 2: Module-level Trade-offs.} Subgraph-extraction is the main efficiency bottleneck for large-scale graphs. Overuse of semantic augmentation (e.g., using LLMs) in subgraph-extraction can compromise downstream path-retrieval quality. Targeting semantic augmentation in the path-retrieval module is a more cost-effective strategy to maintain high quality.
    %For quality, due to the sequential dependency between the two modules, excessive semantic augmentation (e.g., using LLMs) in subgraph-extraction may lead to the loss of critical information in path-retrieval, increasing cost and degrading quality. Conversely, applying semantic augmentation solely to the path-retrieval can enhance quality while minimizing cost.
        
    %\item 
    \textbf{Finding 3: Method-level Trade-offs.} Structure-based methods offer higher efficiency and lower cost, but semantic-augmented methods are essential for quality, especially for complex (two- or multi-hop) queries. 
    One-way semantic methods (OSAR) balance quality and efficiency better than interactive methods (ISAR).
    %Among semantic-augmented methods, one-way methods (OSAR) achieve a better balance between quality and efficiency than interactive methods (ISAR), with comparable cost. 

    %\item 
    \textbf{Finding 4: Model-level Trade-offs.} 
    EEMs are more cost effective in token and GPU usage, while LLMs generally provide better quality, particularly for interactive methods (ISAR). However, LLMs can exhibit unstable generation in some cases, affecting consistency. A balanced approach includes applying LLM-based augmentation in path-retrieval or EEM-based augmentation in both modules.
    %To balance quality and cost, either apply LLM-based semantic augmentation solely to the path-retrieval module or use EEM-based semantic augmentation in both modules.

    %\item 
    \textbf{Finding 5: A Promising yet Underexplored Strategy.} Based on our analysis, the effective approach for GraphRAG would be the one that combines structure based methods for subgraph-extraction with OSAR for path-retrieval. Despite its potential, this strategy remains underexplored in existing works. Also, it opens up opportunities for further optimizations for handling multi-hop and multi-answer queries, as well as improving overall efficiency. 




%\vspace{-5pt}
\subsection{Efficiency  of the Subgraph-Extraction}
\label{subsec:eeSE}

We explore promising directions for faster subgraph-extraction, supported by empirical analysis, including limitations, trade-offs, and research opportunities.\footnote{See technical report B.7 for details.} Below, we briefly outline each direction.

%and discuss their trade-offs, limitations, and future research opportunities. A detailed discussion is provided in our technical report~\ref{Efficiency-SE}; below we briefly outline each direction.

{\bf Computational acceleration} yields significant efficiency improvement for SBE methods (e.g., PPR) through 
%two categories of techniques: 
approximate and distributed algorithms (Table~\ref{tab: PPR-Acceleration}). 
%As shown in Table~\ref{tab: PPR-Acceleration}, representative methods from both categories achieve substantial speedups. 
%However, approximate methods may reduce subgraph recall, while distributed ones rely on specific system configurations.
Approximate methods may lower Recall; distributed solutions depend on system setup.


{\bf Precomputation-based acceleration} pre-generates subgraphs (e.g., via clustering or community detection) to scale down the search space for the extraction process in order to reduce runtime overhead. 
%prepares the graph in advance by generating subgraphs using methods such as clustering or community detection, allowing the extraction process to focus on a smaller portion of the graph. 
As shown in Table~\ref{tab: SE-Acceleration}, our implemented prototype demonstrates the potential of this direction, but it also incurs preprocessing overhead and may limit the ability to access relevant information spread across different subgraphs.



{\bf Vector database-based acceleration} encodes graph components (e.g., nodes or triples) into embedding vectors and stores them in a vector database, thus efficiently identifying query-relevant components based on semantic similarity, followed by subgraph construction centered on them. 
%As shown in Table~\ref{tab: SE-Acceleration}, our prototype shows the potential of this direction, achieving significant speedups with relatively high recall. However, the retrieved components are not necessarily structurally connected, often producing subgraphs with multiple weakly connected components (WCCs), which may limit their effectiveness in multi-hop reasoning
Table~\ref{tab: SE-Acceleration} shows our prototype achieves notable speedups with high Recall. However, the retrieved components often form multiple weakly connected subgraphs, which may hinder multi-hop reasoning.




\begin{table}[t]
\centering
\caption{{\small Computational acceleration of PPR}} % 空标题仅用于占位
\vspace{-10pt}
\resizebox{\linewidth}{!}{%
    \begin{tabular}{l|l|l||l|l}  
        % 左表标题和列名
        \toprule
        \multicolumn{3}{c||}{\textbf{Approximate PPR (Freebase)}} & 
        \multicolumn{2}{c}{\textbf{Distributed PPR}} \\
        \midrule
        \multicolumn{1}{c|}{\textbf{Method}} & 
        \multicolumn{1}{c|}{\textbf{Recall}} & 
        \multicolumn{1}{c||}{\textbf{Ave. Time (s)}} & 
        \multicolumn{1}{c|}{\textbf{Method}} & 
        \multicolumn{1}{c}{\textbf{Speedup}} \\   
        \midrule
        
        % 左表数据 || 右表数据
        RBS~\cite{RBS} & 0.31 & \textbf{0.51} & 
        HGPA~\cite{HGPA} & 3.4--4.1$\times$ \\   
        
        Fora~\cite{Fora} & 0.18 & 7.61 & 
        PAFO~\cite{PAFO} & 58.7$\times$ \\   
        
        TopPPR~\cite{Topppr} & 0.44 & 42.08 & 
        Delta-Push~\cite{Delta-Push} & \textbf{123--162$\times$} \\   
        
        \midrule
        \textbf{Standard PPR} & \textbf{0.96} & 75.29 & 
        \textbf{Standard PPR} & 1$\times$ (baseline) \\   
        \bottomrule
    \end{tabular}%
}
\label{tab: PPR-Acceleration}
\end{table}


\begin{table}[t]  
\vspace{-10pt}
        \centering   
    \caption{{\small Performance and Cost Analysis of the SE Acceleration Methods on Freebase (About 100M Nodes, 300M Edges)}}  
    \vspace{-10pt}
    \label{tab: SE-Acceleration} 
    \resizebox{0.46\textwidth}{!}{%

    \begin{tabular}{l|l|l|l|l|l}  
        \toprule  
        \textbf{Method}                    & \textbf{Pre-time} & \textbf{Online-time} & \textbf{Recall} & \textbf{F1} & \textbf{WCC} \\   
        \midrule         
        Precomputation                    & 19373s                             & 19.02s                               & 0.52          & 0.0021      & \textbf{1}                                      \\
        Vector Database                    & 14570s                             & \textbf{0.85s}                                & 0.65          & 0.0023      & 12.86                                  \\   
        \midrule  
        \textbf{Standard PPR}                           & \textbf{0s}                                 & 75.29s                      & \textbf{0.96} & \textbf{0.0038} & \textbf{1}                            \\
        \bottomrule  
    \end{tabular}  
    }
\end{table}  

\begin{table}[h]  
\vspace{-10pt}
    \centering  
    \small
    \caption{\small Performance of strategies designed to mitigate the limitations of LLM-generated outputs }
    \vspace{-10pt}
    \label{tab:PR-LLM}  
    \resizebox{0.48\textwidth}{!}{%  
    \begin{tabular}{l|l|l|l|l}  
        \toprule  
        \textbf{Method} & \textbf{CWQ} & \textbf{WebQSP} & \textbf{GrailQA} & \textbf{WebQuestions} \\  
        \midrule  
        \textbf{SBE+OSAR-LLMs} & 0.304 & 0.401 & \ul{0.401} & 0.328 \\  \midrule 
        \textbf{SBE+OSAR-LLMs (M-LLMs)} & \ul{0.349} & \textbf{0.438} & 0.381 & \ul{0.348} \\  
        \textbf{SBE+OSAR-LLMs (EEMs-S)} & \textbf{0.381} & \ul{0.427} & \textbf{0.476} & \textbf{0.353} \\  
        \bottomrule  
    \end{tabular}  
    }
\end{table}  

\subsection{Generation Quality of LLMs}
\label{subsec:LLMsQ}



For the issue of LLMs generating fewer or lower-quality components, we take the path-retrieval module as a example and explore two effective strategies: {\bf Multi-round LLMs (M-LLMs)} 
%invoke the LLM multiple times to iteratively refine and improve output quality; 
conduct iterative refinement via repeated LLM calls.
{\bf EEMs supplementation (EEMs-S)} compensates for insufficient outputs by incorporating top-ranked candidate paths from EEMs. % to fill in the missing content.
Table~\ref{tab:PR-LLM} shows both strategies improve performance across most datasets (applied in the SBE+OSAR-LLMs instance), though with added computational overhead of multiple model calls and EEM-based ranking steps.\footnote{More details and analysis are in technical report B.9. }